% concatenated from main.tex
\pdfoutput=1
\documentclass[11pt]{article}

\usepackage[T1]{fontenc}
\usepackage[utf8]{inputenc}
\usepackage[a4paper,margin=1in]{geometry}
\usepackage{lmodern}
\usepackage{amsmath,amssymb,amsfonts,mathtools,amsthm}
\usepackage{array,booktabs}
\usepackage{graphicx}
\usepackage{microtype}
\usepackage{xcolor}
\usepackage{cite}
\usepackage[linesnumbered,lined,ruled,commentsnumbered]{algorithm2e}
\usepackage[colorlinks=true,linkcolor=blue,citecolor=blue,urlcolor=blue]{hyperref}

\hypersetup{
  pdftitle={Data-driven feedback rectification of switched linear systems},
  pdfauthor={Philipp Schmitz, Hannes Gernandt, Maria C. Honecker, Karl Worthmann}
}

\setlength{\emergencystretch}{3em}
\allowdisplaybreaks

\newcommand{\R}{\mathbb{R}}
\newcommand{\C}{\mathbb{C}}
\newcommand{\N}{\mathbb{N}}
\DeclareMathOperator{\im}{im}

\theoremstyle{plain}
\newtheorem{theorem}{Theorem}
\newtheorem{lemma}[theorem]{Lemma}
\newtheorem{proposition}[theorem]{Proposition}
\newtheorem{corollary}[theorem]{Corollary}
\theoremstyle{definition}
\newtheorem{definition}[theorem]{Definition}
\newtheorem{assumption}[theorem]{Assumption}
\theoremstyle{remark}
\newtheorem{remark}[theorem]{Remark}


\title{Data-driven feedback rectification of switched linear systems}

\author{%
\parbox{0.96\textwidth}{%
\centering
Philipp Schmitz$^{1,*}$, Hannes Gernandt$^{2,3}$,\\
Maria C. Honecker$^{4,5}$, and Karl Worthmann$^{1}$\\[0.75em]
\footnotesize
$^{1}$Optimization-based Control Group, Technische Universit\"at Ilmenau, 98693 Ilmenau, Germany\\
$^{2}$University of Wuppertal, Gau\ss stra\ss e 20, 42119 Wuppertal, Germany\\
$^{3}$Fraunhofer Research Institution for Energy Infrastructures and Geotechnologies IEG,\\ 03046 Cottbus, Germany\\
$^{4}$Chair of Intelligent Control Systems, RWTH Aachen University,\\ Kopernikusstra\ss e 16, 52062 Aachen, Germany\\
$^{5}$cigus GmbH, Magirus-Deutz-Stra\ss e 13, 89077 Ulm, Germany\\[0.5em]
$^{*}$Corresponding author: \href{mailto:philipp.schmitz@tu-ilmenau.de}{philipp.schmitz@tu-ilmenau.de}%
}%
}

\date{\today}

\begin{document}
\maketitle

\begin{abstract}
In this paper, a data-driven method for the computation of stabilizing state-feedbacks is proposed that leads to a rectified eigenstructure of switched linear systems. This means that all switching subsystems have the same sets of eigenvectors and the rectification allows to compute a common quadratic Lyapunov function ensuring asymptotic stability of the closed-loop system. The method is illustrated for two examples including a model of an Aerosonde.
\end{abstract}

\noindent\textbf{Keywords:} switched linear systems, Willems' fundamental lemma, eigenstructure assignment, common quadratic Lyapunov functions, feedback, data-driven control design

\section{Introduction}

Switched systems provide a natural modeling framework for dynamical processes whose evolution is governed by different modes of operation. Such switching phenomena occur in a broad range of applications, including power converters in electrical systems~\cite{trenn2012switched},
aircraft control systems~\cite{WangZS2016}, precision motion systems~\cite{HeertjesN2012}, electro-hydraulic systems~\cite{YuanDSB2017}, and coupled mechanical systems~\cite{Sajja2019} to name but a few. Accordingly, switched systems have been investigated extensively during the last decades, resulting in several monographs and a large body of stability and control-theoretic results, see e.g.\ the textbooks~\cite{Liberzon.2003,zhao2017control} or the survey paper~\cite{lin2009stability}.

A central difficulty in the analysis of switched systems is that stability does, in general, not only depend
on the stability properties of the individual subsystems, but also on the switching signal. In particular, even switching between stable subsystems may lead to instability~\cite{Liberzon.2003}.
This motivates the study of stability under arbitrary switching, that is, stability guarantees which are independent of the particular switching sequence.
A standard tool for establishing such guarantees is the construction of common Lyapunov functions~\cite{mason2023universal}. Among these, common quadratic Lyapunov functions (CQLFs) are especially attractive, since they provide verifiable and computationally tractable certificates for stability under arbitrary switching~\cite{andersen2022common}.

With the development of stability criteria for switched systems, increasing attention has been devoted to feedback design. Existing approaches include methods based on linear matrix inequalities \cite{Sun05} and $H_\infty$-control tools \cite{Deaecto11}
as well as model-reference adaptive control~\cite{YuanDSB2017} and switched model-predictive control~\cite{ZhangZB2016}. Other contributions address feedback designs that guarantee stability under broad classes of switching signals while also aiming at additional performance objectives; examples include composite nonlinear feedback designs for switched nonlinear~\cite{LinPB1998} and networked systems~\cite{li2018improved,zhao2022event}.

A particularly constructive line of work relies on eigenstructure assignment. In contrast to purely numerical design approaches, eigenstructure-based methods are closely related to classical linear control techniques and provide direct insight into the closed-loop dynamics.

For linear time-invariant systems with multiple inputs, assigning the
closed-loop poles by state feedback generally leaves freedom in the associated
closed-loop eigenvectors. Moore's seminal result~\cite{Moore76} characterizes
the admissible eigenvectors for prescribed eigenvalues and thereby makes
explicit the design freedom beyond pole placement. Subsequent parametric
formulations have exploited this freedom, for example, to optimize the
robustness and gain of pole-placing feedbacks, including in the presence of
repeated poles~\cite{Schmid14}, and to connect eigenstructure assignment with
geometric-control constructions~\cite{padula2021eigenstructure}.

Related recent work provides explicit parameterizations of generalized
eigenvector chains, which are needed when the prescribed closed-loop
eigenstructure contains nontrivial Jordan chains
\cite{maynard5761099parameterizing}

For switched linear systems under arbitrary switching, eigenstructure-based
feedback design has been used to impose common structural properties on the
closed-loop modes. The approach in~\cite{WulWirSho09} employs invariant-subspace
constructions to stabilize a class of SISO switched systems, whereas
\cite{Hai16} studies the simultaneous triangularization of the closed-loop
matrices together with arbitrary eigenvalue assignment. The feedback
rectification framework of~\cite{HoneckerGWTR24} instead constructs
mode-dependent feedback matrices such that the closed-loop subsystems share a
common eigenvector basis, applying Moore's eigenstructure-assignment result to
the individual modes.

If the assigned eigenvalues are stable, the common eigenvector basis yields a
CQLF and hence stability under arbitrary switching. At the same time, the
remaining freedom in the eigenvector selection can be exploited for additional
performance objectives, as demonstrated for monotone tracking
in~\cite{HoneckerSWR24} and non-overshooting output shaping
in~\cite{WulffHSR26}.

The existing feedback rectification methods are model-based: they require exact knowledge of the subsystem matrices and involve the computation of basis vectors from intersections of eigenvector-assignment subspaces. While analytical descriptions of these intersections are available in certain cases, for example for two subsystems, their computation becomes increasingly involved as the number of modes grows. This motivates the question whether feedback rectification can be carried out directly from measured data, without first relying on an explicit model-based representation of all subsystems.

Data-driven control has become an active research area, in particular following Willems' fundamental lemma~\cite{Willems2005}, which allows the behavior of controllable linear time-invariant systems to be represented from sufficiently rich input--state or input-output data, see also the review papers~\cite{markovsky2021behavioral,Faulwasser23}.
This perspective has led to numerous data-driven stabilization and control methods for linear systems; see, for example,~\cite{van2020willems} and data-driven pole assignment~\cite{celi23}. More recently, data-driven methods have also been applied to switched linear systems in discrete time~\cite{eising2024data}, in continuous time \cite{bianchi2025data}, for output feedback control~\cite{hu2023data}, and also to nonlinear systems~\cite{Ju2026}.

These approaches typically rely on data informativity conditions to obtain stabilizing controllers from data, possibly without knowing the switching signal during the data collection phase. The setting considered in the present paper is different and complementary: our goal is not only to stabilize the switched system, but to construct mode-dependent feedbacks that rectify the closed-loop eigenstructure. This provides an explicit CQLF and allows the designer to prescribe desired closed-loop eigenvalue locations.

In this paper, we propose a data-driven method for the rectification and stabilization of switched linear systems. We work in discrete time and assume that, during an offline phase, sufficiently informative input--state data are collected for each subsystem. Based on these data, we derive algebraic conditions that characterize the feasibility of a common closed-loop eigenvector basis and provide explicit formulas for the corresponding rectifying feedback matrices. The resulting closed-loop systems share the same eigenvectors, while their eigenvalues may be chosen mode-dependently inside the unit disk. Consequently, a CQLF can be constructed directly from the data, yielding asymptotic stability under arbitrary switching, provided that the active mode is available for the implementation of the corresponding feedback.

The main contributions of this paper are summarized as follows.
\begin{itemize}
    \item We formulate feedback rectification for discrete-time switched linear systems for multiple eigenvalues and relate the existence of stabilizing rectifying feedback matrices to a common closed-loop eigenstructure.
    \item Using input--state data and Willems' fundamental lemma, we derive a data-based representation of the eigenvector assignment subspaces that appear in Moore's eigenstructure assignment theorem.
    \item We provide a data-driven characterization of the parametric subspace conditions required for feedback rectification.
    \item We obtain an explicit data-driven formula for the mode-dependent feedback matrices.
    \item We show how the rectified eigenstructure yields a CQLF directly from data and summarize the resulting design in an algorithm.
\end{itemize}

The paper is organized as follows. Section~\ref{sec:prelim} recalls preliminaries on switched systems, CQLFs, Moore's eigenstructure assignment theorem, and model-based feedback rectification. Section~\ref{sec:main_data} contains the main data-driven results, including the data-based rectifiability condition, the feedback formula, and the resulting algorithm for constructing stabilizing rectifying feedbacks and a CQLF. Section~\ref{sec:example} illustrates the proposed method by numerical examples. Finally, Section~\ref{sec:conclusion} concludes the paper and outlines directions for future work.

\section{Preliminaries}
\label{sec:prelim}
In this section, we recall stability, stabilization and existence conditions for CQLFs in Subsection~\ref{ssec:cqlfs}. Subsection~\ref{ssec:feedbackrectification} presents the feedback rectification together with a~sufficient condition based on  \cite{HoneckerGWTR24}.
\subsection{Asymptotic stability, stabilization, and CQLFs}
\label{ssec:cqlfs}
We consider a switched linear system that is composed by the constituent discrete-time linear time-invariant (LTI) systems given by
\begin{equation}
    x_{k+1} = A_{q} x_k + B_{q} u_k,\quad k\geq 0,
\end{equation}
with $A_q\in\R^{n\times n}$, $B_q\in\R^{n\times m}$ and $q\in\mathcal Q=\{1,2,\dots, \overline q\}$,  where $k\in\N$. These subsystems  constitute the discrete-time switched system
\begin{equation}
    \label{eq:switched_sys}
    x_{k+1} = A_{\sigma(k)} x_k + B_{\sigma(k)} u_k,\quad k\geq 0,
\end{equation}
with switching signal $\sigma: \N \to \mathcal Q$. For the stability analysis we consider the homogeneous system
\begin{align}
    \label{eq:homo}
 x_{k+1} = A_{\sigma(k)}x_{k},\quad  x_0=x^0, \quad k\geq 0,
\end{align}
and call the equilibrium $x^*=0$ \emph{asymptotically stable} if $\lim_{k\rightarrow\infty} x_k=x^*=0$. A \emph{common quadratic Lyapunov function} (CQLF) is a function $\mathcal{V}(x)=\tfrac{1}{2}x^\top Qx$ for some positive definite $Q\in\mathbb{R}^{n\times n}$ which is decreasing along trajectories, i.e.\ there exists $\alpha>0$ such that $\mathcal{V}(x_{k+1})-\mathcal{V}(x_{k})\leq -\alpha \|x_k\|^2$ for all $k\geq 0$.

The aim of this paper is to design based on input--state data a feedback law of the form
\begin{equation}
    \label{eq:feedback_law}
    u_k = F_{\sigma(k)} x_k
\end{equation} with $F_{\sigma(k)}\in \R^{m\times n}$ which asymptotically stabilizes the closed-loop switched system
\begin{equation}
    x_{k+1} = (A_{\sigma(k)} + B_{\sigma(k)}F_{\sigma(k)}) x_k,\quad k\geq 0,
\end{equation}
under arbitrary but known switching signal $\sigma$, by constructing a CQLF for the closed-loop system.

In the following proposition,
we recall a result on the exponential stability of the homogeneous system \eqref{eq:homo} using common quadratic Lyapunov functions. To this end, we use the \emph{spectral radius} $\rho(A)$ of a matrix $A\in\mathbb{R}^{n\times n}$ which is the maximum modulus of its eigenvalues. The proof of the proposition is presented in Section~\ref{sec:prop_lyapunov} in the appendix.
\begin{proposition}
\label{prop:lyapunov}
Given the homogeneous switched system~\eqref{eq:homo} with $\rho(A_q)<1$ for all $q\in\mathcal Q$. If there exists a regular $T\in\C^{n\times n}$ such that $T^{-1}A_q T$ is upper triangular for all $q\in\mathcal Q$, then there exists a common quadratic Lyapunov function. Consequently, the origin of~\eqref{eq:homo} is exponentially stable under arbitrary switching. If in addition $T^{-1}A_q T$ are diagonal for all $q\in\mathcal Q$, then the Lyapunov function is given by
    \[
    \mathcal{V}(x)=\tfrac{1}{2}x^\top T^{-*}T^{-1}x\quad \text{for all $x\in\mathbb{R}^n$}.
    \]
\end{proposition}

\subsection{Model-based feedback rectification}
\label{ssec:feedbackrectification}
The design of the switched feedback law~\eqref{eq:feedback_law} relies on the seminal result of Moore \cite{Moore76}, which characterizes the eigenvalue assignment for a single LTI system
\begin{equation}
    \label{eq:LTI_sys}
    x_{k+1} = A x_k + B u_k
\end{equation}
with respect to the closed-loop matrix $A + BF$. For $\lambda\in\C$ consider the matrices $N(\lambda)$ and $M(\lambda)$ satisfying
\begin{equation}
    \label{eq:NlambMlamb}
    \operatorname{ker}\begin{bmatrix}
        \lambda I - A & B
    \end{bmatrix} = \operatorname{im}\begin{bmatrix}
        N(\lambda)\\ M(\lambda)
    \end{bmatrix}.
\end{equation}

The following result is an extension of \cite{Moore76} from \cite{KleinMoore77} which characterizes the assignment of multiple eigenvalues.
\begin{theorem}
\label{thm:moore}
Given $(A,B)\in \mathbb R^{n\times n}\times \mathbb R^{n\times m}$ and the self-conjugate set $\Lambda=\{\lambda_1,\dots,\lambda_n\}$. Then there exists $F\in \mathbb R^{m\times n}$ such that $(A+BF)v_i = \lambda_i v_i$, $v_i\neq 0$, for all $i\in\{1,\dots, n\}$ if and only if the following three conditions are satisfied:
\begin{itemize}
    \item[\rm (i)] $\{v_1,\dots, v_n\}\subset \mathbb C^{n}$ is linearly independent;
    \item[\rm (ii)] there exists an involution, i.e.\ a self-inverse mapping
$\pi:\{1,\ldots,n\}\to\{1,\ldots,n\}$ such that
\[
    \lambda_{\pi(i)}=\overline{\lambda_i},
    \qquad
    v_{\pi(i)}=\overline{v_i},
    \qquad i=1,\ldots,n;
\]
    \item[\rm (iii)] $v_i\in \operatorname{im} N(\lambda_i)$ for all $i=1,\ldots,n$,
\end{itemize}
where $\bar v$ and $\bar \lambda$ denote the complex conjugate of $v$ and $\lambda$, respectively.
\end{theorem}

If the conditions of Theorem~\ref{thm:moore} are satisfied, the feedback matrix $F$ can be computed explicitly. Specifically, for eigenvectors $v_i = N(\lambda_i) w_i$ with suitable $w_i$, the feedback matrix $F$ satisfies
\begin{equation}
    \label{eq:F}
    F\begin{bmatrix}
        v_1 & \dots & v_n
    \end{bmatrix} = -\begin{bmatrix}
        M(\lambda_1)w_1 & \dots & M(\lambda_n)w_n
        \end{bmatrix}.
\end{equation}
Since $\begin{bmatrix}
        v_1 & \dots & v_n
    \end{bmatrix}$ is invertible we can rewrite
\begin{align}
\label{eq:F_explicit_with_inverse}
F=-\begin{bmatrix}
        M(\lambda_1)w_1 &\!\! \dots &\!\! M(\lambda_n)w_n
        \end{bmatrix}\begin{bmatrix}
        v_1 &\!\! \dots &\!\! v_n
    \end{bmatrix}^{-1}.
\end{align}
\begin{remark}
\label{rem:real_feedback}
Note that the matrix $F$ in \eqref{eq:F_explicit_with_inverse} has indeed real entries, i.e.\ $F\in\mathbb{R}^{m\times n}$, as stated in Theorem~\ref{thm:moore} which can be seen from separating the columns into real and imaginary parts. Here, without loss of generality we assume that the eigenvectors are sorted in such a way that the first $2n_1$ eigenvectors correspond to the non-real eigenvalues and that they appear pair-wise with a vector and its complex-conjugate, i.e.\ $v_{2i}=\overline{v_{2i+1}}$ for $i=1,\ldots n_1$ and that the remaining $n-2n_1$ eigenvectors correspond to the real eigenvalues, i.e.\ $v_{i}=\overline{v_{i}}$ for $i=2n_1+1,\ldots,n$. With this assumption and using the abbreviation $v=v^R+\mathrm{i}v^I$ for the decomposition into real and imaginary parts, i.e.\ $v^R, v^I\in\mathbb{R}^n$, and writing $\hat{w}_i:= M(\lambda_i)w_i$ for $i=1,\ldots,n$, we obtain
\begin{align*}
&~~~~\begin{bmatrix}
        M(\lambda_1)w_1 &\!\! \dots &\!\! M(\lambda_n)w_n
        \end{bmatrix}\begin{bmatrix}
        v_1 &\!\! \dots &\!\! v_n
    \end{bmatrix}^{-1}\\
    & = \left[\begin{smallmatrix} \hat{w}_1
^R+ \mathrm{i}\hat w_1^I &\hat{w}_1
^R- \mathrm{i}\hat w_1^I  & \ldots & \hat{w}_{n_1}
^R- \mathrm{i}\hat w_{n_1}^I &\hat{w}_{2n_{1}+1}&\ldots & \hat w_n
        \end{smallmatrix}\right]\\ & ~~~\cdot\begin{bmatrix}
        v_1^R+{\rm i} v_1^I &  v_1^R-{\rm i} v_1^I  &\!\! \dots & v_{2n_1+1}& \ldots&\!\! v_n
    \end{bmatrix}^{-1} \\
    &= \begin{bmatrix} \hat w_1^R& \hat w_1^I & \ldots & \hat w_{n_1}^R& \hat w_{n_1}^I & \hat w_{2n_1+1} \ldots &\!\!  \hat w_n
        \end{bmatrix} \\ &~~~~\left[\begin{smallmatrix}
\mathrm{diag}\left(\begin{bmatrix}
            1 & 1 \\ \mathrm{i} & -\mathrm{i}
        \end{bmatrix}\right)   &0\\0& I_{n-2n_1}      \end{smallmatrix} \right]\left[\begin{smallmatrix}
\mathrm{diag}\left(\begin{bmatrix}
            1 & 1 \\ \mathrm{i} & -\mathrm{i}
        \end{bmatrix}\right)   &0\\0& I_{n-2n_1}      \end{smallmatrix}\right]^{-1}\\
        &~~~~\begin{bmatrix}
        v_1^R & v_1^I &\!\! \dots &\!\! v_{n_1}^R & v_{n_1}^I& v_{2n_1+1} & \ldots & v_n
    \end{bmatrix}^{-1}
    \\&=\begin{bmatrix} \hat w_1^R& \hat w_1^I \, \ldots \,\hat w_{n_1}^R& \hat w_{n_1}^I & \hat w_{2n_1+1} \,\ldots \, \hat w_n
        \end{bmatrix}\\
        &~~~~\begin{bmatrix} v_1^R& v_1^I \, \ldots \, v_{n_1}^R& v_{n_1}^I & v_{2n_1+1}\,  \ldots \,  v_n
        \end{bmatrix}^{-1} \in\mathbb{R}^{m\times n}.
\end{align*}
From this calculation it can be seen that the feedback matrix $F$ still has real entries when the eigenvectors are not sorted pairwise by applying suitable perturbation matrices.
\end{remark}

We recall the definition of feedback rectification from \cite{HoneckerGWTR24} slightly extended to more than two switching subsystems.
\begin{definition}
\label{def:fbrect}
Let $A_q\in\R^{n\times n}$ and $B_q\in\R^{n\times m}$, $q\in\mathcal{Q}$, and  let $D\subset\C^{\bar q}$.
Then the set of matrices $\{(A_q,B_q)\}_{q\in \mathcal{Q}}$ is called \emph{feedback rectifiable over $D$} if there exists $F_q\in\R^{m\times n}$, tuples of complex numbers $(\lambda_{i}^{(1)},\ldots,\lambda_{i}^{(\bar q)})\in D$, $i=1,\ldots,n$, and a linearly independent set $\{v_1,\ldots,v_n\}$ in $\C^n$ such that the following holds
\begin{align}
\label{def:ferec}
(A_q+B_qF_q)v_i=\lambda_i^{(q)}v_i,
\end{align}
for all $i=1,\ldots,n$ and all $q=1,\ldots,\bar q$.
\end{definition}

A sufficient condition for feedback rectifiability over $D$ is given by the following result, which is a slight generalization of \cite{HoneckerGWTR24} to switched systems with more than two subsystems and directly follows from Theorem~\ref{thm:moore}. Similarly to a single LTI system case, we consider for $\lambda\in\C$ and $q\in\mathcal Q$ matrices $N_q(\lambda)$, $M_q(\lambda)$ satisfying
\begin{equation}
    \ker \begin{bmatrix}
        \lambda I-A_q & B_q
    \end{bmatrix} = \im \begin{bmatrix}
        N_q(\lambda)\\
        M_q(\lambda)
    \end{bmatrix}.
\end{equation}

\begin{proposition}
\label{prop:rect}
Let $A_q\in\R^{n\times n}$ and $B_q\in\R^{n\times m}$, $q\in\mathcal{Q}$, and let $D\subset\C^{\bar q}$.
Then the family $\{(A_q,B_q)\}_{q\in\mathcal Q}$ is feedback
rectifiable over $D$ if there exist indexed eigenvalue tuples
\[
    \bigl(
        \lambda_i^{(1)},\ldots,\lambda_i^{(\bar q)}
    \bigr)
    \in D,
    \qquad i=1,\ldots,n,
\]
an involution
$\pi:\{1,\ldots,n\}\to\{1,\ldots,n\}$, $\pi(\pi(i))=i$,
and linearly independent vectors
$v_1,\ldots,v_n\in\mathbb C^n$ such that, for every
$i=1,\ldots,n$ we have  $v_i\in
        \bigcap_{q=1}^{\bar q}
        \operatorname{im}N_q\bigl(\lambda_i^{(q)}\bigr)$ and the following condition
\begin{align}
    \lambda_{\pi(i)}^{(q)}
        &= \overline{\lambda_i^{(q)}}
        \quad \text{for every } q\in\mathcal Q, \quad \text{and} \quad    v_{\pi(i)}
        = \overline{v_i}. \label{eq:tuple_conjugacy}
\end{align}
\end{proposition}

\begin{remark}
In the case of two switching systems, i.e.\ $\bar q=|\mathcal{Q}|=2$ one can use row-reduced echelon forms over matrices of rational functions in two variables to compute the subspace intersection in Proposition~\ref{prop:rect} explicitly \cite{HoneckerGWTR24}. In principle, this can be extended to more than two subsystems, although this might lead to complex expressions of rational functions in the subspace representation.
\end{remark}

\section{Data-driven feedback rectification}
\label{sec:main_data}

In this section, we present our main results on data-driven feedback rectification. First, we recall in Subsection~\ref{ssec:data_LTI} the fundamental lemma by Willems et al. \cite{Willems2005}. In Subsection~\ref{ssec:data_main} we present the main result where we obtain a data-driven condition to verify the sufficient condition for rectifiability from Proposition~\ref{prop:rect} and a formula for computation of a rectifying feedback.

\subsection{Single LTI systems: Data-driven feedback rectification}
\label{ssec:data_LTI}

Suppose that we have measured an input--state trajectory $(x,u)$ of the LTI system~\eqref{eq:LTI_sys}, and collect these data in matrices
\begin{equation}
    \label{eq:data_matrices}
    \begin{aligned}
    U &= \begin{bmatrix}
        u_0 & \dots & u_{N-2}
    \end{bmatrix},\quad X = \begin{bmatrix}
        x_0 & \dots & x_{N-2}
    \end{bmatrix},\\
    X_+ &= \begin{bmatrix}
        x_1 & \dots & x_{N-1}
    \end{bmatrix}.
    \end{aligned}
\end{equation}
In this case, the data matrices satisfy
\begin{equation}
\label{eq:data_dynamcis}
    X_+ = A X + BU.
\end{equation}
A necessary condition on the data, in order to completely describe the system, is
\begin{equation}
    \label{eq:rank}
    \operatorname{rank}\begin{bmatrix}
        X\\ U
    \end{bmatrix} = n+m.
\end{equation}
In this case the system matrices $A$ and $B$ can be directly reconstructed from the data,
\begin{equation}
    \begin{bmatrix}
        A & B
    \end{bmatrix} = X_+ \begin{bmatrix}
        X\\U
    \end{bmatrix}^\dagger,
\end{equation}
where $M^\dagger$ denotes the Moore--Penrose pseudo inverse of a matrix $M$. A sufficient input-design condition for the data to satisfy \eqref{eq:rank} follows e.g.\ by the Willems et al.'s fundamental lemma \cite[Theorem~1]{Willems2005}, see also \cite[Thorem~1]{van2020willems}.
\begin{lemma}
\label{lem:FL}
    Let $(A,B)$ be controllable and consider an input--state $(u,x)$ trajectory of system~\eqref{eq:LTI_sys} of length $N$. If $u$ is persistently exciting of order $L=1+n$, that is, the Hankel matrix
    \begin{equation}
        \begin{bmatrix}
            u_0 & \dots & u_{N-L}\\
            \vdots & \ddots & \vdots\\
            u_{L-1} & \dots & u_{N-1}
        \end{bmatrix}\in\mathbb{R}^{mL\times (N-L+1)}
    \end{equation}
    has full row rank, then \eqref{eq:rank} holds.
\end{lemma}

If no single consecutive data trajectory satisfying~\eqref{eq:rank} can be measured, for instance due to interrupted measurements, one may instead rely on multiple fragmented datasets. As shown in \cite[Theorem~2]{van2020willems}, such datasets can collectively provide a complete description of the LTI system.
\begin{lemma}
\label{lem:FLfrag}
    Let $(A,B)$ be controllable and consider input--state trajectories $(u^i, x^i)$ of length $N_i$, $i=1,\dots,K$. If the inputs $(u^i, \dots, u^K)$ are collectively persistently exciting of the order $L=1+n$, that is the mosaic-Hankel matrix
    \begin{equation}
        \begin{bmatrix}
            u_0^1 & \dots & u_{N_1-L}^1 & \dots & u_{0}^K & \dots & u_{N_K-L}^K\\
            \vdots & \ddots & \vdots & & \vdots & \ddots & \vdots\\
            u_{L-1}^1 & \dots & u_{N_1-1}^1 & \dots & u_{L-1}^K & \dots & u_{N_K-1}^K
        \end{bmatrix}
    \end{equation}
    has full row rank, then the data matrices
    \begin{equation*}
    \begin{aligned}
        U &= \begin{bmatrix}
            u_0^1&\dots & u_{N_1-2}^1 & \dots & u_{0}^K & \dots & u_{N_K-2}^K
        \end{bmatrix},\\[0.5em]
        X &= \begin{bmatrix}
            x_0^1&\dots & x_{N_1-2}^1 & \dots & x_{0}^K & \dots & x_{N_K-2}^K
        \end{bmatrix},\\[0.5em]
        X_+ &= \begin{bmatrix}
            x_1^1&\dots & x_{N_1-1}^1 & \dots & x_{1}^K & \dots & x_{N_K-1}^K
        \end{bmatrix}
    \end{aligned}
    \end{equation*}
    satisfy \eqref{eq:data_dynamcis} and \eqref{eq:rank}.
\end{lemma}
Given sufficiently informative input--state data, $\operatorname{ker}\begin{bmatrix}
    \lambda I-A & B
\end{bmatrix}$ and the relation \eqref{eq:NlambMlamb} can be reformulated in a data-based manner.
\begin{lemma}
    \label{lem:data_kernel}
    Consider the data matrices $U,X,X_+$ in~\eqref{eq:data_matrices} and suppose \eqref{eq:rank}. For every $\lambda\in\mathbb C$ the following holds:
    \begin{equation}
    \label{eq:data_kernel}
        \operatorname{ker}\begin{bmatrix}
            \lambda I - A & B
        \end{bmatrix} = \operatorname{im}\begin{bmatrix} N(\lambda)\\ M(\lambda) \end{bmatrix}=  \begin{bmatrix} X\\ -U\end{bmatrix}\operatorname{ker}\begin{bmatrix}
        \lambda X-X_+\end{bmatrix}.
        \end{equation}
\end{lemma}
\begin{proof}
    For any pair of compatible matrices $P$ and $Q$, where $Q$ is surjective, it holds that $\operatorname{ker}(P) = Q\operatorname{ker}(PQ)$. Let $P = \begin{bmatrix}
        \lambda I - A & B
    \end{bmatrix}$ and $Q = [\begin{smallmatrix}X\\ -U\end{smallmatrix}]$. By \eqref{eq:rank}, $Q$ is surjective.
    Moreover, \eqref{eq:data_dynamcis} yields
    \begin{equation}
        \lambda X - X_+ = \begin{bmatrix}
            \lambda I - A & B
        \end{bmatrix}\begin{bmatrix}X\\ -U\end{bmatrix} = PQ.
    \end{equation}
    This yields the claim.
\end{proof}
Consequently, this allows Moore's theorem to be restated in terms of data, see also \cite{celi23} for a data-driven pole-placement result.
\begin{theorem}
\label{thm:moore_dd}
    Consider the data matrices $U,X,X_+$ in~\eqref{eq:data_matrices} and suppose \eqref{eq:rank}.
    Given the self-conjugate set $\Lambda=\{\lambda_1,\dots,\lambda_n\}$, there exists $F\in \mathbb R^{m\times n}$ such that $(A+BF)v_i = \lambda_i v_i$, $v_i\neq 0$, for all $i\in\{1,\dots, n\}$ if and only if the following three conditions are satisfied:
\begin{itemize}
    \item[\rm (i)] $\{v_1,\dots, v_n\}\subset \mathbb C^{n}$ is linearly independent;
  \item[\rm (ii)] there exists an involution
$\pi:\{1,\ldots,n\}\to\{1,\ldots,n\}$ such that, for every
$i=1,\ldots,n$,
\[
    \lambda_{\pi(i)}=\overline{\lambda_i},
    \qquad
    v_{\pi(i)}=\overline{v_i};
\]
    \item[\rm (iii)] $v_i\in X \operatorname{ker}(\lambda_i X-X_+)$ for all $i=1,\ldots,n$.
\end{itemize}

Let $g_i\in \ker(\lambda_i X-X_+)$ such that $v_i=Xg_i$ for $i=1,\dots, n$. Then the feedback matrix $F\in\mathbb{R}^{m\times n}$ is given by
\begin{equation}
\label{eq:ddF}
    F = U\begin{bmatrix}
        g_1 & \dots &  g_n
    \end{bmatrix} \begin{bmatrix}
        v_1 & \dots & v_n
    \end{bmatrix}^{-1}.
\end{equation}
\end{theorem}

\begin{proof}
    The claimed equivalence follows directly from Theorem~\ref{thm:moore} and Lemma~\ref{lem:data_kernel}. We show formula~\eqref{eq:ddF}. Let $w_i$, $i=1,\dots, n$ such that
    \begin{equation}
        v_i = X g_i = N(\lambda_i)w_i,\quad -U g_i = M(\lambda_i) w_i
    \end{equation}
    which exist by \eqref{eq:data_kernel} in Lemma~\ref{lem:data_kernel}. Together with~\eqref{eq:F} we find~\eqref{eq:ddF}.
\end{proof}

\subsection{Switched LTI systems: Data-driven feedback rectification}
\label{ssec:data_main}
Now we focus on the switched system~\eqref{eq:switched_sys}. In an offline phase, we collect for each LTI subsystem $(A_q,B_q)$, $q\in\mathcal Q$, input--state data $(u^q, x^q)$, and define data matrices
\begin{equation}
\label{eq:data_matq}
\begin{aligned}
    U^q &= \begin{bmatrix}
        u_0^q&\dots &u_{N_q-2}^q
    \end{bmatrix},\quad X^q = \begin{bmatrix}
        x_0^q & \dots & x_{N_q-2}^q
    \end{bmatrix}\\
    X_+^q &= \begin{bmatrix}
        x_1^q&\dots & x_{N_q-1}^q
    \end{bmatrix}.
\end{aligned}
\end{equation}
We assume that the data is sufficiently informative and satisfies for every $q\in\mathcal Q$
\begin{equation}
    \label{eq:rankq}
    \operatorname{rank}\begin{bmatrix}
        X^q\\ U^q
    \end{bmatrix} = n + m.
\end{equation}
As stated in Lemma~\ref{lem:FL}, provided that $(A^q,B^q)$ is controllable, condition~\eqref{eq:rankq} is guaranteed if the input sequence $u^q$ is persistently exciting of order $n+1$. In the case that the excitation for the LTI subsystems is interrupted before sufficiently long consecutive trajectories can be recorded, e.g.\ by switching events, one may instead collect multiple shorter fragments for each LTI system which are collectively persistently exciting of order $n+1$, see Lemma~\ref{lem:FLfrag}.

Next we characterize the intersection of the images of the matrices $N_q(\lambda_q)$, $q\in\mathcal Q$, in terms of data.
\begin{lemma}
    Consider data matrices $U^q, X^q, X_+^q$ in \eqref{eq:data_matq} such that \eqref{eq:rankq} holds for all $q\in\mathcal Q$. Let $\lambda_1,\dots, \lambda_{\bar q}\in\mathbb C$ and define
\begin{equation}
\label{eq:curlyX}
\!\!\!    \mathcal X(\lambda_1,\dots, \lambda_{\bar q}) \!:=\!
    \begin{bmatrix}
        X_+^1 \!- \!\lambda_1 X^1 & 0 \\
        0&\!\!\!\!\! \mathrm{diag}(X_+^i \!- \!\lambda_i X^i)_{i=2}^{\bar q}\\
        \begin{bmatrix}
            X^1\\ \vdots \\ X^1
        \end{bmatrix} &\!\!\!\!\! - \mathrm{diag}(X^i)_{i=2}^{\bar q}
    \end{bmatrix}\!.
\end{equation}
Then
\begin{equation}
    \bigcap_{q\in\mathcal Q} \operatorname{im}N_q(\lambda_q) = \begin{bmatrix} X^1 &\!\! 0 &\!\! \dots &\!\! 0\end{bmatrix} \operatorname{ker}\mathcal X(\lambda_1,\dots, \lambda_{\bar q}).
\end{equation}
\end{lemma}
\begin{proof}
    Let $v\in \begin{bmatrix}
        X^1 & 0 &\dots & 0
    \end{bmatrix} \ker \mathcal X(\lambda_1,\dots, \lambda_{\bar q})$ and consider a vector $g = \begin{bmatrix}
        g_{1}^\top & \dots & g_{\bar q}^\top
    \end{bmatrix}{}^\top\in \ker \mathcal X(\lambda_1,\dots, \lambda_{\bar q})$ such that $v = \begin{bmatrix}
        X^1 & 0 &\dots & 0
    \end{bmatrix} g$. Then $g_q\in \ker(X_+^q-\lambda_q X^q)$ and $v = X^1 g_1 = X^q g_q$ for all $q\in\mathcal Q$. By Lemma~\ref{lem:data_kernel}, $v = X^q g_q \in \operatorname{im} N_q(\lambda_q)$ for all $q\in\mathcal Q$.

    Let $v\in \bigcap_{q\in\mathcal Q} \operatorname{im}N_q(\lambda_q)$. Then by Lemma~\ref{lem:data_kernel} there exist $g_q\in \ker(X_+^q-\lambda_q X^q)$ such that $v = X^q g_q$. Therefore, $g = \begin{bmatrix}
        g_{1}^\top & \dots & g_{\bar q}^\top
    \end{bmatrix}{}^\top\in \ker \mathcal X(\lambda_1,\dots, \lambda_{\bar q})$ and $v\in \begin{bmatrix}
        X^1 & 0 &\dots & 0
    \end{bmatrix} \ker \mathcal X(\lambda_1,\dots, \lambda_{\bar q})$.
\end{proof}

Note that in \eqref{eq:curlyX}, $X^1$ plays the role of a reference matrix. This choice is arbitrary and any of the matrices $X^1,\dots, X^{\bar q}$ may serve this purpose, leaving the space $\begin{bmatrix} X^1 & 0 & \dots & 0\end{bmatrix} \operatorname{ker}\mathcal X(\lambda_1,\dots, \lambda_{\bar q})$ invariant.

\begin{theorem}
\label{thm:data_rectification}
    Consider data matrices $U^q, X^q, X_+^q$ in \eqref{eq:data_matq} such that \eqref{eq:rankq} holds for all $q\in\mathcal Q$. Given self-conjugate sets $\Lambda_q = \{\lambda_{1}^{(q)},\dots, \lambda_{n}^{(q)}\}$ for $q\in\mathcal Q$, consider the matrices $\mathcal X(\lambda_{i}^{(1)},\dots, \lambda_{i}^{(\bar q)})$ defined as in \eqref{eq:curlyX}. The following conditions are satisfied:
    \begin{itemize}
        \item[\rm (i)] $\{v_1,\dots, v_n\}\subset \C^n$ is linearly independent,
        \item[\rm (ii)] there exists an involution
$\pi:\{1,\ldots,n\}\to\{1,\ldots,n\}$, $\pi(\pi(i))=i$ such that conjugacy condition \eqref{eq:tuple_conjugacy} holds;
        \item[\rm (iii)] $v_i\in \begin{bmatrix}
            X^1 & 0 & \dots & 0
        \end{bmatrix} \ker \mathcal X(\lambda_{i}^{(1)},\dots, \lambda_{i}^{(\bar q)})$,
    \end{itemize}
if and only if for every $q\in\mathcal Q$ there exists $F_q\in\R^{m\times n}$ such that $(A_q+B_qF_q)v_i = \lambda_{i}^{(q)} v_i$ for all $i=1,\dots,n$.

    Let $g_i=\begin{bmatrix}
        g_{i,1}^\top&\dots&g_{i,\bar q}^\top
    \end{bmatrix}{}^\top\in \ker \mathcal X(\lambda_{i}^{(1)},\dots, \lambda_{i}^{(\bar q)})$ such that $v_i=X^1g_{i,1}$ for $i=1,\dots, n$. Then the feedback matrix $F_q\in\mathbb{R}^{m\times n}$, $q\in\mathcal Q$, is given by
    \begin{equation}
    \label{eq:F_data}
    F_q = U^q \begin{bmatrix}
        g_{1,q} & \dots & g_{n,q}
    \end{bmatrix}
        \begin{bmatrix}
            v_1 & \dots & v_n
        \end{bmatrix}^{-1}.
    \end{equation}
\end{theorem}
\begin{proof}
    The equivalence of the rectifiability of $\{(A_q, B_q)\}_{q\in\mathcal Q}$ and the statements (i)--(iii) follows immediately from Theorem~\ref{thm:moore} and Lemma~\ref{lem:data_kernel}. The representation~\eqref{eq:F_data} for the feedback matrices $F_1,\dots, F_{\bar q}$ follows with \eqref{eq:ddF} in Theorem~\ref{thm:moore_dd} together with Lemma~\ref{lem:data_kernel}. That $F_q$ are real-valued for all $q\in\mathcal{Q}$ can be seen as in Remark~\ref{rem:real_feedback}.
\end{proof}

The procedure to calculate the feedback matrices based on data is summarized in Algorithm~\ref{alg:1}.

\begin{algorithm}
\caption{Data-driven feedback rectification}\label{alg:1}
\KwData{$\{(X^q,X_+^q,U^q)\}_{q\in\mathcal Q}$ satisfying \eqref{eq:data_matq}, \eqref{eq:rankq}}
\KwResult{matrices $\{F_q\}_{q\in \mathcal Q}$, CQLF: $x\mapsto \tfrac{1}{2} x^\top P x$}
choose eigenvalue candidates $\{\lambda_i^{(q)}\}$ according to Theorem~\ref{thm:data_rectification} with $\lvert \lambda_i^{(q)}\rvert<1$ \;
find $g_i = \begin{bmatrix}
    g_{i,1}^\top & \dots & g_{i,\bar q}^\top
\end{bmatrix}{}^\top \in \ker\mathcal X(\lambda_i^{(1)},\dots, \lambda_i^{(\bar q)})$ such that
$V=\begin{bmatrix}
    X^1 g_{1,1} & \dots & X^1 g_{n,1}
\end{bmatrix}$ is invertible\;
$F_q \gets U^q \begin{bmatrix}
    g_{1,q} & \dots & g_{n,q}
\end{bmatrix} V^{-1}$ \;
$P \gets V^{-*} V^{-1}$\;
\end{algorithm}

\begin{remark}
Although we assume the switching signal to be known, the proposed scheme can be combined with the data-driven mode detection approach in~\cite{eising2024data}. If the CQLF value ceases to decrease, this may indicate an unobserved mode switch. Hence, mode detection can be triggered to identify the active subsystem and select the corresponding feedback gain. This extension requires additional data-informativity assumptions and a detection procedure that is sufficiently fast relative to the system dynamics. Under these conditions, the approach can also be applied when the switching signal is not known a priori.
\end{remark}

\subsection{Nonlinear program formulation}
The computation of kernel elements in Algorithm~\ref{alg:1} comes with the additional constraint that the assembled matrix $V$ is invertible, which is difficult to enforce by standard linear-algebraic techniques. To address this issue, we propose a heuristic approach based on solving the following nonlinear optimization problem:
\begin{subequations}
\label{eq:opt}
    \begin{equation}
        \label{eq:opta}
        \operatorname*{minimize}_{\{\Lambda,\ V,\ g\}}\ \lVert VV^\top -I\rVert_\mathrm{F}^2\quad \text{s.t.}
    \end{equation}
    \begin{align}
        \label{eq:optb}
        \Lambda &= \{\lambda_i^{(q)}\}_{i=1,\dots, n,\ q\in\mathcal Q} \subset [-\alpha, \alpha]\\
        V &= \begin{bmatrix}
            v_1 & \dots & v_n
        \end{bmatrix}\in\mathbb R^{n\times n},\\
        g &= \begin{bmatrix}
            g_{1,1} &\!\!\dots\!\!& g_{1,\bar q} &\!\! \dots\!\! & g_{n,1} &\! \!\dots\!\! & g_{n,\bar q}
        \end{bmatrix}\in\mathbb R^{(N-1)\bar q n}\\
        \label{eq:opte}
        0&=c_{i,q}(\Lambda,V,g):=\begin{bmatrix}
            v_i\\
            \lambda_i^{(q)} v_i
        \end{bmatrix} - \begin{bmatrix}
            X^q\\
            X_+^q
        \end{bmatrix}g_{i,q}.
    \end{align}
\end{subequations}

Observe that the objective function in \eqref{eq:opta} is quartic, while constraint~\eqref{eq:opte} is quadratic, rendering the optimization problem~\eqref{eq:opt} nonconvex. The user-defined parameter $\alpha\in(0,1)$ in \eqref{eq:opt} ensures that the constructed feedback stabilizes the closed-loop system. A small value of $\lVert VV^\top -I\rVert_\mathrm{F}$ promotes invertibility of $V$ and drives it towards orthogonality. Indeed,
\begin{equation}
\label{eq:bd}
    1>\lVert VV^\top -I\rVert_\mathrm{F}^2 = \sum_{i=1}^n (\sigma_i^2 - 1)^2
\end{equation} implies that all singular values $\sigma_i$ of $V$ are positive and $V$ is invertible. Note that the bound~\eqref{eq:bd} is conservative and invertibility of $V$ might hold for considerably larger values of $\lVert VV^\top -I\rVert_\mathrm{F}^2$.

To solve the constrained optimization problem~\eqref{eq:opt}, we apply the augmented Lagrangian method \cite{nocedal2006numerical}. Given $\mu\in\R^{2n^2\bar q}$ and $\nu>0$ the augmented Lagrangian reads
\begin{equation*}
\begin{split}
    &\mathcal L(\Lambda,\ V,\ g,\mu;\nu):=\\
    &\frac{1}{2}\lVert VV^\top -I\rVert_\mathrm{F}^2 + \mu^\top c(\Lambda,\ V,\ g) + \frac{\nu}{2} \lVert c(\Lambda,\ V,\ g)\rVert_2^2,
\end{split}
\end{equation*}
where
\begin{equation}
    c(\Lambda,\ V,\ g) := \begin{bmatrix}
        c_{1,1}(\Lambda,V, g)&
        \ldots &
        c_{n,\bar q}(\Lambda, V, g)
    \end{bmatrix}^\top.
\end{equation}
Then the minimization problem
\begin{equation}
    \operatorname*{minimize}_{\{\Lambda,\ V,\ g\}}\ \mathcal L(\Lambda,V,g,\mu_k;\nu_k)
\end{equation}
subject to the bound constraint~\eqref{eq:optb} is solved successively via the limited-memory BFGS algorithm, see, e.g., \cite{nocedal2006numerical}. The penalty factors $\nu_k>0$ form a nondecreasing sequence chosen adaptively with $\nu_k\rightarrow\infty$ and the Lagrangian multipliers $\mu_k$ are updated appropriately between iterations. The solution of each subproblem is used as initialization for the next iteration, such that the resulting sequence of iterates approximates a (local) minimizer of nonlinear program~\eqref{eq:opt}.

\section{Example}
\label{sec:example}
We demonstrate the functionality of the proposed algorithm with two examples.
\subsection{Example with randomly chosen matrix entries}
We consider a switched system with $\bar q=4$ switching modes, input dimension $m=2$ and state dimension $n=6$. The subsystems are randomly generated while ensuring feedback rectifiability. All numerical values following are rounded to two significant digits. The system matrices are given by
\begin{align*}
    A_{1} &= \left[\begin{smallmatrix}
    0.06 & -0.69 & 0.43 & 0.21 & -0.16 & 0.24 \\
    0.28 & -0.71 & -0.35 & -0.64 & -0.35 & 0.24 \\
    0.20 & -0.01 & -0.36 & -0.18 & 1.26 & -0.70 \\
    0.71 & -0.07 & 0.11 & 0.37 & -0.49 & 0.19 \\
    -0.02 & 0.09 & -0.15 & -0.08 & -0.44 & -0.58 \\
    -0.25 & -0.08 & -0.09 & -0.21 & 0.06 & 0.80
\end{smallmatrix}\right]\\
    A_{2} &= \left[\begin{smallmatrix}
    0.64 & -0.34 & 0.61 & 0.11 & -0.04 & 0.06 \\
    0.49 & -0.47 & -0.21 & -0.45 & -0.19 & -0.07 \\
    0.34 & -0.11 & 0.16 & -0.32 & 0.53 & -0.41 \\
    0.20 & -0.40 & -0.48 & 0.20 & -0.21 & 0.45 \\
    0.16 & 0.08 & -0.24 & -0.47 & -0.47 & -0.33 \\
    0.04 & 0 & -0.17 & -0.12 & 0.17 & 0.77
\end{smallmatrix}\right]\\
    A_{3} &= \left[\begin{smallmatrix}
    0.58 & -0.19 & 0.32 & 0.62 & 0.03 & -0.12 \\
    0.67 & -0.25 & -0.68 & -0.49 & 0.09 & -0.06 \\
    0.52 & -0.15 & -0.30 & -0.09 & 0.83 & -0.33 \\
    0.64 & -0.09 & -0.42 & -0.10 & -0.24 & 0.51 \\
    0.18 & -0.08 & -0.16 & -0.11 & -0.55 & -0.33 \\
    0.02 & 0.03 & -0.03 & 0 & 0.04 & 0.70
\end{smallmatrix}\right]\\
 A_{4} &= \left[\begin{smallmatrix}
    1.07 & -0.65 & 0.61 & 0.58 & 0.16 & 0.08 \\
    0.24 & -0.47 & -0.50 & -0.56 & -0.09 & -0.02 \\
    -0.23 & 0.03 & -0.17 & 0 & 0.49 & -0.55 \\
    0.03 & -0.52 & 0.12 & -0.12 & -0.66 & 0.48 \\
    -0.50 & 0.10 & -0.09 & -0.05 & -0.86 & -0.53 \\
    -0.60 & 0.08 & -0.10 & -0.08 & -0.16 & 0.65
\end{smallmatrix}\right]
\end{align*}
\begin{equation*}
\begin{aligned}
 B_{1} &= \left[\begin{smallmatrix}
    -0.35 & 0.10 \\
    -0.38 & -0.14 \\
    0.46 & 0.50 \\
    0.27 & -0.19 \\
    0.19 & 0.21 \\
    -0.11 & 0.14
\end{smallmatrix}\right]&
    B_{2} &= \left[\begin{smallmatrix}
    -0.11 & -0.42 \\
    0.26 & -0.07 \\
    -0.03 & -0.35 \\
    -0.32 & 0.41 \\
    -0.46 & -0.27 \\
    -0.21 & -0.01
\end{smallmatrix}\right]\\
    B_{3} &= \left[\begin{smallmatrix}
    -0.16 & 0.34 \\
    -0.02 & 0.27 \\
    0.35 & 0 \\
    0.41 & 0.09 \\
    0.35 & -0.16 \\
    0 & 0.03
\end{smallmatrix}\right]&
    B_{4} &= \left[\begin{smallmatrix}
    -0.49 & -0.22 \\
    0.20 & 0.13 \\
    0.48 & 0.12 \\
    -0.02 & 0.26 \\
    0.40 & 0.22 \\
    0.46 & 0.28
\end{smallmatrix}\right].
\end{aligned}
\end{equation*}
Since we consider a data-driven setting, each subsystem is simulated to produce an input--state trajectory $(u^q,x^q)$, from which the corresponding data matrices $U^q,X^q,X_+^q$ as in \eqref{eq:data_matq} are constructed. Given these data matrices, the intersection matrix in \eqref{eq:curlyX} can be computed. By applying Algorithm \ref{alg:1} we find a switched state feedback with matrices
\begin{align*}
    F_{1} &= \left[\begin{smallmatrix}
    -0.98 & -0.58 & 0.05 & -0.67 & -0.67 & 0.67 \\
    0.98 & 0.11 & 0.68 & 0.98 & -0.72 & -0.10
\end{smallmatrix}\right] \\
    F_{2} &= \left[\begin{smallmatrix}
    0.17 & -0.01 & -0.83 & -0.51 & 0.69 & 0.28 \\
    0.30 & 0.34 & 0.53 & -0.88 & -0.27 & 0.08
\end{smallmatrix}\right] \\
    F_{3} &= \left[\begin{smallmatrix}
    -0.79 & -0.20 & 0.83 & 0.26 & -0.64 & -0.32 \\
    -0.62 & -0.95 & 0.85 & -0.10 & -0.38 & 0.20
\end{smallmatrix}\right] \\
    F_{4} &= \left[\begin{smallmatrix}
    0.73 & -0.77 & 0.46 & -0.12 & 0.11 & 0.31 \\
    0.94 & 0.97 & -0.42 & 0.47 & 0.50 & -0.31
\end{smallmatrix}\right]
\end{align*}
The common eigenvectors assigned form the columns of
\[
V = \left[\begin{smallmatrix}
    -0.08 & -0.10 & -0.15 & -0.76 & -0.09 & 0.30 \\
    -0.53 & -0.63 & -0.15 & -0.15 & -0.09 & 0.06 \\
    0.25 & -0.50 & -0.10 & -0.02 & 0.16 & 0.56 \\
    -0.51 & 0.10 & 0.09 & -0.40 & -0.14 & -0.33 \\
    -0.45 & 0.38 & 0 & 0.07 & 0 & 0.17 \\
    0.01 & -0.01 & 0 & -0.20 & 0 & -0.52
\end{smallmatrix}\right]
\]
and the corresponding Lyapunov matrix reads
\[
P = \left[\begin{smallmatrix}
    8.95 & 7.04 & -15.31 & -17.18 & 0 & 0 \\
    7.04 & 9.48 & -15.90 & -17.84 & 0 & 0 \\
    -15.31 & -15.90 & 41.00 & 27.76 & 14.67 & 19.96 \\
    -17.18 & -17.84 & 27.76 & 47.41 & -13.07 & -17.79 \\
    0 & 0 & 14.67 & -13.07 & 26.71 & 33.39 \\
    0 & 0 & 19.96 & -17.79 & 33.39 & 47.63
\end{smallmatrix}\right].
\]
The eigenvalues of the closed-loop subsystems can be seen in Tab.~\ref{tab:ev_rand}.
\begin{table}
\caption{Eigenvalues of the closed-loop matrices $A_q+B_q F_q$ of the random switched system.}
\label{tab:ev_rand}
\centering
\begin{tabular}{rrrrrrr}
\toprule
$q$ & $\lambda_1$ & $\lambda_2$ & $\lambda_3$ & $\lambda_4$ & $\lambda_5$ & $\lambda_6$ \\
\midrule
$1$ & $-0.7$ & $-0.7$ & $0.0$ & $0.0$ & $0.7$ & $0.7$ \\
$2$ & $-0.7$ & $-0.7$ & $0.0$ & $0.0$ & $0.7$ & $0.7$ \\
$3$ & $-0.7$ & $-0.7$ & $0.0$ & $0.0$ & $0.7$ & $0.7$ \\
$4$ & $-0.7$ & $-0.7$ & $0.0$ & $0.0$ & $0.7$ & $0.7$ \\
\bottomrule
\end{tabular}
\end{table}
Figure~\ref{fig:rand} depicts the Euclidean norm along a closed-loop state trajectory together with the values of the Lyapunov function for known random switching events.
\begin{figure}[htbp!]
    \centering
    \includegraphics[width=\linewidth]{rnf.pdf}
    \caption{Random switched system. Euclidean norm and CQLF along state trajectory, under random switching signal.}
    \label{fig:rand}
\end{figure}

\subsection{Unmanned Aerial Vehicle (UAV)}
In \cite{honecker2025development} an Unmanned Aerial Vehicle (UAV) is modeled with $\bar q = 2$ switching modes, input dimension $m=5$ and state dimension $n=7$. Mode one is the nominal mode and mode two models a defective mode which models mainly a loss of yaw controllability. This example is chosen as  in \cite{honecker2025development} it is demonstrated to be difficult to solve with \cite{HoneckerGWTR24} and restrictions on the eigenvalue pairs apply. All numerical values following are rounded to three significant digits. The system matrices of the individual subsystems are given by
\begin{align*}
    A_{1} &= \left[\begin{smallmatrix}
    0.998 & -0.003 & -0.012 & -0.490 & 0 & 0 & 0 \\
    -0.005 & 0.004 & 0 & 0.002 & 0 & 0 & 0 \\
    -0.038 & -0.241 & 0.966 & 0.009 & 0 & 0 & 0 \\
    -0.001 & -0.010 & 0.049 & 1.000 & 0 & 0 & 0 \\
    0 & 0 & 0 & 0 & 0.957 & 0 & 0 \\
    0 & 0 & 0 & 0 & -0.137 & 0.450 & 0.226 \\
    0 & 0 & 0 & 0 & 0.190 & 0.017 & 0.605
\end{smallmatrix}\right] \\
    A_{2} &= \left[\begin{smallmatrix}
    0.998 & -0.005 & -0.012 & -0.490 & 0 & 0 & 0 \\
    -0.005 & 0.004 & 0 & 0.002 & 0 & 0 & 0 \\
    -0.043 & 0 & 0.966 & 0.011 & 0 & 0 & 0 \\
    -0.001 & 0 & 0.049 & 1.000 & 0 & 0 & 0 \\
    0 & 0 & 0 & 0 & 0.957 & 0 & 0 \\
    0 & 0 & 0 & 0 & -0.137 & 0.450 & 0.226 \\
    0 & 0 & 0 & 0 & 0.190 & 0.017 & 0.605
\end{smallmatrix}\right] \\
   B_{1} &= \left[\begin{smallmatrix}
    0.005 & 0 & 0 & 0.002 & -0.001 \\
    0.104 & 0 & 0 & 0 & 0.009 \\
    -1.871 & 0 & 0 & 0.047 & 0.293 \\
    -0.047 & 0 & 0 & 0.001 & 0.007 \\
    0 & 0 & -0.008 & 0 & 0 \\
    0 & 4.136 & 5.258 & 0 & 0 \\
    0 & 0.110 & -0.802 & 0 & 0
\end{smallmatrix}\right] \\
    B_{2} &= \left[\begin{smallmatrix}
    0.005 & 0 & 0 & 0.003 & -0.001 \\
    0.104 & 0 & 0 & 0 & 0 \\
    -1.757 & 0 & 0 & 0 & 0.303 \\
    -0.044 & 0 & 0 & 0 & 0.008 \\
    0.054 & 0 & 0 & 0 & 0 \\
    -0.005 & 4.136 & 0 & 0 & 0 \\
    0.006 & 0.110 & 0 & 0 & 0
\end{smallmatrix}\right].
\end{align*}

The calculated feedback matrices are given by
\begin{align*}
    F_{1} &= \left[\begin{smallmatrix}
    -0.929 & 3.879 & 1.066 & 7.476 & -91.866 & -1.393 & -0.085 \\
    -0.003 & -0.045 & -0.039 & -0.211 & -80.031 & -0.138 & -0.443 \\
    0.001 & 0.002 & 0 & -0.002 & 55.243 & -0.002 & 0.331 \\
    -94.054 & 4.166 & 43.979 & 460.423 & -710.992 & -10.396 & -0.818 \\
    9.642 & 25.039 & -6.277 & -54.316 & -491.222 & -7.511 & -0.433
\end{smallmatrix}\right] \\
    F_{2} &= \left[\begin{smallmatrix}
    0 & 0.002 & 0 & 0 & -6.527 & 0 & -0.001 \\
    -0.001 & -0.001 & -0.045 & -0.244 & -9.981 & -0.114 & -0.065 \\
    0 & 0 & 0 & 0 & 0 & 0 & 0 \\
    -93.900 & -1.599 & 42.906 & 454.823 & -379.027 & -5.876 & -0.578 \\
    0.440 & 0.004 & -5.793 & -26.959 & -59.410 & -0.327 & -0.030
\end{smallmatrix}\right]
\end{align*}
with eigenvector matrix
\[
V = \left[\begin{smallmatrix}
    0.004 & -0.035 & 0.680 & 0 & 0.046 & 0.738 & -0.003 \\
    0.007 & -0.234 & -0.003 & -0.004 & 0 & -0.004 & -0.972 \\
    0 & 0.135 & -0.484 & 0 & -0.706 & 0.498 & -0.034 \\
    0 & -0.003 & 0.101 & 0 & 0.002 & -0.085 & 0.001 \\
    -0.006 & 0 & 0 & 0.004 & 0 & 0 & 0 \\
    0.365 & -0.839 & -0.018 & -0.303 & -0.167 & -0.014 & 0.206 \\
    0.648 & 0.007 & -0.002 & 0.762 & -0.002 & -0.002 & -0.002
\end{smallmatrix}\right]
\]
and Lyapunov matrix given by
\[
P = \left[\begin{smallmatrix}
    0.996 & -0.006 & -0.074 & -0.777 & 0.727 & 0.004 & 0.002 \\
    -0.006 & 1.002 & -0.003 & -0.009 & 5.356 & 0.018 & 0.005 \\
    -0.074 & -0.003 & 1.990 & 10.623 & -1.987 & -0.028 & -0.007 \\
    -0.777 & -0.009 & 10.623 & 114.044 & -20.669 & -0.300 & -0.041 \\
    0.727 & 5.356 & -1.987 & -20.669 & 25288.313 & 85.901 & 30.171 \\
    0.004 & 0.018 & -0.028 & -0.300 & 85.901 & 1.292 & 0.103 \\
    0.002 & 0.005 & -0.007 & -0.041 & 30.171 & 0.103 & 1.036
\end{smallmatrix}\right].
\]
The eigenvalues of the corresponding closed-loop subsystems are shown in Tab.~\ref{tab:ev_uav}. Fig.~\ref{fig:aero} depicts the Euclidean norm and the Lyapunov function along some state trajectory, given a known random switching signal.

\begin{figure}[htbp!]
    \centering
    \includegraphics[width=\linewidth]{aeroprobe.pdf}
    \caption{UAV. Euclidean norm and CQLF along state trajectory, under random switching signal.}
    \label{fig:aero}
\end{figure}
\begin{table} \caption{Numerical eigenvalues of the closed-loop matrices $A_q+B_qF_q$ for the UAV switched system.} \label{tab:ev_uav} \centering {\setlength{\tabcolsep}{3pt} \begin{tabular}{@{}rrrrrrrr@{}} \toprule $q$ & $\lambda_1$ & $\lambda_2$ & $\lambda_3$ & $\lambda_4$ & $\lambda_5$ & $\lambda_6$ & $\lambda_7$ \\ \midrule $1$ & $-0.800$ & $-0.155$ & $0.031$ & $0.694$ & $0.744$ & $0.787$ & $0.800$ \\ $2$ & $-0.800$ & $0.004$ & $0.005$ & $0.591$ & $0.609$ & $0.744$ & $0.787$ \\ \bottomrule \end{tabular} } \end{table}

\section{Conclusion and outlook}
\label{sec:conclusion}
In this paper, we present a data-driven method for the stabilization of switched linear systems. Using input--state data collected during an offline phase, we derive state feedbacks that rectify the eigenstructure of switched linear systems.
As a consequence, this allows us to obtain a common quadratic Lyapunov function purely based on the data.
Future work will include the incorporation of additive noise in the data using data-informativity and including a data-driven mode detection for the stabilization with an unknown switching signal.

\section*{Acknowledgments}
Philipp Schmitz gratefully acknowledges funding by the Carl Zeiss Foundation (VERNEDCT, Project No.~2021-10-003). Hannes Gernandt was funded by the Deutsche Forschungsgemeinschaft (DFG, German Research Foundation), Project-ID~531152215, CRC~1701. Karl Worthmann gratefully acknowledges support by the DFG, Project-ID~554600805.

% The original bibliography database was not included in the uploaded source.
% Once refs.bib is added, this branch automatically produces the bibliography.
\IfFileExists{refs.bib}{%
  \bibliographystyle{IEEEtran}
  \bibliography{refs}
}{%
  \section*{References}
  \noindent\textbf{Source asset required before submission.}
  The bibliography database \texttt{refs.bib} was not included with the uploaded manuscript.
  Add the original file next to \texttt{main.tex}; the bibliography will then be generated automatically.
}

\appendix
\section{Proof of Proposition~\ref{prop:lyapunov}}
\label{sec:prop_lyapunov}
\begin{proof}
Since $T^{-1}A_qT$, $q\in\mathcal Q$, is upper triangular, we write $T^{-1}A_qT=\Lambda_q+N_q$ for $\Lambda_q$ diagonal containing the eigenvalues and
$N_q$ is strictly upper triangular satisfying
\[
\|\Lambda_q\|_2=\rho(A_q)<1.
\]
Now introduce, for $\varepsilon\in (0,1]$, a diagonal scaling
$D_\varepsilon:=\operatorname{diag}(\varepsilon^k)_{k=0}^{n-1}\in\mathbb{R}^{n\times n}$.

Then
\[
\widetilde U_q(\varepsilon):=D_\varepsilon^{-1}(T^{-1}A_qT)D_\varepsilon
=
\Lambda_q+D_\varepsilon^{-1}N_qD_\varepsilon.
\]
Since $N_q$ is strictly upper triangular, for $i<j$,
\[
\bigl(D_\varepsilon^{-1}N_qD_\varepsilon\bigr)_{ij}
=
\varepsilon^{\,j-i}(N_q)_{ij}.
\]
Thus, every off-diagonal entry is multiplied by a factor
$\varepsilon^{j-i}$, hence by at most $\varepsilon$.
Therefore,
\[
\|D_\varepsilon^{-1}N_qD_\varepsilon\|_F
\le \varepsilon\,\|N_q\|_F.
\]
Let
$c:=\max_{q\in\mathcal Q} \|N_q\|_F$. Then, for every $q\in\mathcal Q$,
\[
\|D_\varepsilon^{-1}N_qD_\varepsilon\|_2
\le
\|D_\varepsilon^{-1}N_qD_\varepsilon\|_F
\le \varepsilon c.
\]
Hence
\[
\|\widetilde U_q(\varepsilon)\|_2
\le
\|\Lambda_q\|_2+\|D_\varepsilon^{-1}N_qD_\varepsilon\|_2
\le \max_{q\in\mathcal Q}\rho(A_q)+\varepsilon c.
\]
Choose a sufficiently small $\varepsilon\in(0,1]$ such that
$\alpha:=\max_{q\in\mathcal Q}\rho(A_q)+\varepsilon c<1$. Then
\[
\|\widetilde U_q(\varepsilon)\|_2\le \alpha<1
\qquad\forall q\in\mathcal Q.
\]
Now define
$S:=TD_\varepsilon$, then
\[
S^{-1}A_qS
=
D_\varepsilon^{-1}T^{-1}A_qTD_\varepsilon
=
\widetilde U_q(\varepsilon),
\]
implies
\[
\|S^{-1}A_qS\|_2\le \alpha<1
\qquad\forall q\in\mathcal Q.
\]
Equivalently,

\[
A_q^\top P A_q \preceq \alpha^2 P,
\qquad
P:=S^{-*}S^{-1}
=
T^{-*}D_\varepsilon^{-\top}D_\varepsilon^{-1}T^{-1}.
\]
Thus, $\mathcal{V}(x)=x^\top Px$ is a common quadratic Lyapunov function. The proof of the exponential stability is straightforward. If the matrices are jointly diagonalizable, then $N_q=0$ for all $q\in\mathcal{Q}$ and therefore we obtain $c=0$ in the above estimates and we can choose $\varepsilon=1$ to fulfill the estimates.
\end{proof}
\end{document}
